%% ma: L12-B14-0061
%% lop: 12
%% chuong: 5
%% bai: 14
%% dang: 12.14.21
%% loai: TN4
%% muc: NB
%% dap_an: "C"
%% nguon: "To_n_12_PT_M_t_Ph_ng.pdf (Đề kiểm tra đánh giá kiến thức, Chuyên đề Phương trình mặt phẳng, Nguyễn Anh Khôi - Lê Minh Khôi), Câu 16"
%% kien_thuc: "Bài 14 (phương trình mặt phẳng theo đoạn chắn)"
%% trang_thai: cho-duyet
%% ly_do: "Dạng 12.14.21 mới đề xuất, chờ giáo viên duyệt."
%% ghi_chu: "Tài liệu gốc không kèm đáp án; đáp án do Claude giải và kiểm bằng sympy."
\begin{ex}
Trong không gian $Oxyz$, cho $A(1;0;0)$, $B(0;b;0)$, $C(0;0;c)$ với $bc\ne0$. Khi đó phương trình mặt phẳng $(ABC)$ là
\choice{$\dfrac{x}{1}+\dfrac{y}{b}+\dfrac{z}{c}=0$}{$\dfrac{x}{1}-\dfrac{y}{b}-\dfrac{z}{c}=1$}{\True $\dfrac{x}{1}+\dfrac{y}{b}+\dfrac{z}{c}=1$}{$\dfrac{x}{1}-\dfrac{y}{b}-\dfrac{z}{c}=0$}
\loigiai{Phương trình theo đoạn chắn: $(ABC)\colon \dfrac{x}{1}+\dfrac{y}{b}+\dfrac{z}{c}=1$. Chọn C.}
\end{ex}
